% !TeX root = ../main.tex

\chapter{交通知识图谱与场景视觉问答数据集构建}

\section{引言}

\section{多源数据集构建总体框架}

\subsection{交通知识问答与视觉问答的数据需求}

\subsection{两类数据集的功能分工}

\subsection{任务边界与数据构建原则}

\section{面向知识推理的交通事件知识图谱数据集构建}

\subsection{原始数据来源与样本筛选}

\subsection{交通事件实体、属性与关系模式}

\subsection{知识三元组与可追溯证据组织}

\subsection{自然答案生成与知识增强问答样本构建}

\subsection{数据划分与质量控制}

\section{交通场景视觉问答数据集构建}

\subsection{交通场景关系理解与知识推理任务定义}

\subsection{图像、视频与场景样本筛选}

\subsection{交通参与者、场景属性与关系体系}

\subsection{问题类型与问答样本构建}

\subsection{视觉区域、知识依据与答案标注}

\subsection{图文关联样本与不完备关系设置}

\subsection{数据划分与标注一致性检验}

\section{两类问答数据的语义关联与任务接口}

\subsection{交通概念与问题类型的统一组织}

\subsection{知识问答与视觉问答的任务映射}

\subsection{认知数据与四维重建数据的适用边界}

\section{本章小结}
